% Converted from drafts/paper_en.md on 2026-10-21. \method is defined in main.tex. Citation keys are placeholders to be added to main.bib.
\section{Related Work}
\label{sec:related}

\textbf{Subject-driven image generation.} Preserving the identity of a user-provided subject in generated images has been addressed by per-subject optimization \cite{ruiz2023dreambooth, gal2022textual}, by encoder-based adapters that inject image features through cross-attention \cite{ye2023ipadapter, li2023blipdiffusion}, and more recently by treating the reference as additional tokens of a diffusion transformer \cite{tan2024ominicontrol, xiao2024omnigen, wu2025uno}. The token-based formulation scales to large paired datasets whose references deliberately differ from their targets in pose, scene, and attributes \cite{tan2024ominicontrol, xiao2024omnigen}. Our work inherits this data regime and this injection mechanism, but applies them to a \emph{video} backbone.

\textbf{Subject-to-video generation.} Early identity-preserving video methods focused on faces, using identity embeddings as conditions \cite{he2024idanimator, yuan2024consisid}. Open-domain S2V systems condition a video backbone on arbitrary subjects, either by concatenating reference latents along the sequence \cite{liu2025phantom, fei2025skyreelsa2}, by a unified conditioning interface \cite{jiang2025vace}, by multi-modal identity encoders \cite{hu2025hunyuancustom, deng2025magref}, or by a lightweight reference branch that shares the backbone's self-attention \cite{xue2025standin}. All of these are trained on reference--video pairs. OpenS2V-Nexus \cite{yuan2025opens2v} highlights the resulting copy-paste tendency and evaluates naturalness alongside identity, and MuSS \cite{muss2025} proposes a copy-paste rate to quantify it. We adopt Stand-In's reference branch \cite{xue2025standin} as the appearance pathway, but replace video supervision with image supervision and add the mechanisms required to keep motion intact under that supervision.

\textbf{Customizing video models from images.} AnimateDiff \cite{guo2023animatediff} showed that a motion module trained once can be attached to image models customized on image data. Still-Moving \cite{chefer2024stillmoving} addressed the converse problem: adapting a video model's spatial layers on images while keeping its temporal layers, using a motion adapter trained on frozen videos to absorb the static bias of image training and then removed at inference. Our Stage I generalizes this idea to a DiT whose attention is jointly spatiotemporal, and extends it with a signed training regime so that the adapter scale becomes a continuous motion control rather than an on/off switch. Unlike Still-Moving, which adapts the backbone itself, we keep the backbone frozen and use the motion adapter as a fixed carrier on top of which a separate appearance branch is trained.

\textbf{Decoupled conditioning and guidance.} Classifier-free guidance \cite{ho2022cfg} contrasts conditional and unconditional predictions; multi-condition variants apply separate guidance weights to separate conditions \cite{brooks2023instructpix2pix}. Autoguidance \cite{karras2024autoguidance} contrasts a model with a weaker version of itself to improve quality without strengthening the condition. Our paired-difference readout uses a learned null reference not only for guidance at inference but as a \emph{training-time} construction that isolates the reference-dependent component of the prediction. Representation alignment \cite{yu2024repa} regularizes diffusion features toward a frozen visual encoder; we apply it to the reference tokens only.
